% Created 2024-12-10 mar 08:05
% Intended LaTeX compiler: pdflatex
\documentclass[11pt]{article}
\usepackage[utf8]{inputenc}
\usepackage[T1]{fontenc}
\usepackage{graphicx}
\usepackage{longtable}
\usepackage{wrapfig}
\usepackage{rotating}
\usepackage[normalem]{ulem}
\usepackage{amsmath}
\usepackage{amssymb}
\usepackage{capt-of}
\usepackage{hyperref}
\usepackage{minted}
\usepackage[spanish]{inputenc}
\author{Eduardo Alcaraz}
\date{\today}
\title{Proyecto Final}
\hypersetup{
 pdfauthor={Eduardo Alcaraz},
 pdftitle={Proyecto Final},
 pdfkeywords={},
 pdfsubject={},
 pdfcreator={Emacs 28.2 (Org mode 9.5.5)}, 
 pdflang={Spanish}}
\begin{document}

\maketitle
\tableofcontents

El flujo óptico es el patrón del movimiento aparente de los objetos,
superficies y bordes en una escena causado por el movimiento relativo
entre un observador (un ojo o una cámara) y la escena.2​3​ El concepto
de flujo óptico se estudió por primera vez en la década de 1940 y,
finalmente, fue publicado por el psicólogo estadounidense James
J. Gibson4​ como parte de su teoría de la affordance (una acción que un
individuo puede potencialmente realizar en su ambiente). Las
aplicaciones del flujo óptico tales como la detección de movimiento,
la segmentación de objetos, el tiempo hasta la colisión y el enfoque
de cálculo de expansiones, la codificación del movimiento compensado y
la medición de la disparidad estereoscópica utilizan este movimiento
de las superficies y bordes de los objetos.

Actividades

\begin{enumerate}
\item Desarrollar un programa que utilice el flujo óptico, para aplicar a
una imagen o una primitiva de dibujo, las transformaciones
geométricas vistas en clase, (Traslación, Escalamiento, Rotación)
\end{enumerate}


\begin{enumerate}
\item Desarrollar un programa que utilice opengl, para modelar un entorno
completo tipo mundo minecraft de alguna ciudad, paisaje, bosque,
etc. Donde se deben de incluir por lo menos 15 modelos diferentes,
diseñados con primitivas de opengl de igual manera se debe de mover
todo el entorno modelado mediante el flujo óptico aplicando las
transformaciones geométricas vistas en clase, (Traslación,
Escalamiento, Rotación), para esta actividad pueden hacer equipos
máximo de 5 personas.
\end{enumerate}
\end{document}